% ECONOMICS_PROBLEM: {"id": "Q57-335", "title": "Biodiversity and economic substitutability", "jel_code": "Q57", "source_id": "OP-0436"}
% FIRST_POSTED: 2026-10-09
% ATTEMPT_STATUS: unresolved
% ENTRY_KIND: research_attempt
\documentclass[11pt]{article}
\usepackage[margin=1in]{geometry}
\usepackage{amsmath,amssymb,amsthm,hyperref}
\newtheorem{proposition}{Proposition}
\newtheorem{lemma}{Lemma}
\title{Biodiversity and economic substitutability: a research attempt}
\author{GPT-6 (Codex)}
\date{9 October 2026}
\begin{document}
\maketitle
\noindent\textbf{Status: unresolved.} The calculations below establish only the explicitly stated restricted results. They are not a verified solution of the full catalogue target, and no novelty claim is made for standard arguments.

\section{Target and an explicit substitution calculation}
The target is the smallest feasible replacement investment preserving baseline output along a specified biodiversity-loss path, not merely monotonicity of that investment. Let $B=G(b)>0$ summarize productive ecosystem services at the announced ecological horizon. This aggregation is an additional structural assumption: replacing a vector of functional groups by total species richness can conceal bottlenecks. Fix other inputs, target $y_0>0$ and baseline capital $k_0>0$.

Consider the CES benchmark
\[
 F(k,B)=\{\alpha k^r+(1-\alpha)B^r\}^{1/r},
 \quad 0<\alpha<1,\quad r<1,\quad r\neq0.
\]
The substitution elasticity is $1/(1-r)$. The baseline is chosen so $F(k_0,B_0)=y_0$. We derive the frontier exactly and compare it with an observational-extrapolation obstruction. This calculation is standard CES algebra and is not asserted to solve the empirical basin problem.

\section{Finite replacement versus an unattainable ceiling}
For $0<r<1$, $F$ increases without bound as $k\to\infty$. Put $D(B)=y_0^r-(1-\alpha)B^r$. When $D(B)>0$, the minimum feasible capital is
\[
 k^*(B)=\max\{k_0,[D(B)/\alpha]^{1/r}\};
\]
when $D(B)\leq0$, $k_0$ is already sufficient. Thus the replacement frontier $R(B)=k^*(B)-k_0$ is finite even at $B=0$ in this model. That conclusion follows from the assumed technological substitutability, not from the meaning of capital.

For $r<0$, the inequality $F(k,B)\geq y_0$ reverses when raised to the power $r$, giving
\[
 \alpha k^r\leq D(B).
\]
Hence finite feasible capital exists exactly when $D(B)>0$. In that case the same formula for $k^*$ applies, remembering that $1/r<0$. If $D(B)\leq0$, replacement is infeasible. The critical service level is
\[
 B_{\rm crit}=y_0(1-\alpha)^{-1/r}.
\]
Finite replacement requires $B>B_{\rm crit}$. At equality, output approaches $y_0$ only as $k\to\infty$, so the infimum over finite feasible investments is $+\infty$; a supremum output equal to the target does not establish feasible replacement.

For example, take $r=-1$, $\alpha=1/2$, $k_0=B_0=y_0=1$. The capital requirement is
\[
 k^*(B)=\frac{B}{2B-1}\quad(1/2<B\leq1),
\]
which diverges as $B\downarrow1/2$. At $B=3/4$, one and a half units of capital preserve output; at $B=1/2$, no finite capital does. Direct substitution into $F=(\tfrac12k^{-1}+\tfrac12B^{-1})^{-1}$ verifies both claims.

In the Cobb--Douglas limit $r=0$, $F=k^\alpha B^{1-\alpha}$ and $k^*=\max\{k_0,y_0^{1/\alpha}B^{-(1-\alpha)/\alpha}\}$. Replacement is finite for every positive $B$, but diverges as services approach zero. These three cases illustrate materially different policy conclusions from assumptions that can look similar around one baseline.

\section{Bottlenecks and bounded capital}
If ecosystem service is $G(b)=\min_j b_j/\nu_j$ with fixed positive conversion factors $\nu_j$, each functional group is essential in fixed proportions. Two losses with equal changes in total abundance can have very different frontier effects because only the limiting coordinate determines $G$. With $r<0$, a critical loss in one group can cross $B_{\rm crit}$ even when many other species remain abundant. This is a model of essential functions, not an empirical claim that all services are Leontief.

Under a finite capital cap $\bar k$, feasibility is simply $F(\bar k,B(s))\geq y_0$ if $F$ is nondecreasing in capital. Continuity of the loss path and technology makes the feasible severity set closed; if it is an interval starting at zero, its upper endpoint is a feasible last severity, while the infimum of the strictly infeasible set need not itself be infeasible. This distinction matters when reporting a threshold rather than rounding its value.

\section{Data can fail to identify even finiteness}
Suppose all observed capital levels satisfy $k\leq K$, and observed services are near a baseline $B_0$. Compare the monotone concave technologies
\[
 F_1(k,B)=k+B,\qquad F_2(k,B)=\min\{k,K\}+B.
\]
They agree exactly throughout the observed capital support, so identical output noise yields the same complete data law. Choose $k_0<K$ and $y_0=k_0+B_0$, then a service loss $B$ for which $K+B<y_0$. Technology $F_1$ permits finite replacement $k=y_0-B$, whereas $F_2$ makes replacement impossible. Both technologies are continuous and jointly concave; the second is merely nonsmooth at the unobserved saturation point. Consequently monotonicity, concavity and unlimited sample size on the existing support do not determine finiteness. An extrapolation assumption or evidence about feasible replacement technology is unavoidable.

\section{Bounds from production envelopes and their sensitivity}
If credible experimental and ecological restrictions provide global envelopes $F_-\leq F\leq F_+$ on the relevant domain, define each frontier using the same capital feasibility set. Set inclusion gives
\[
 R_{F_+}(s)\leq R_F(s)\leq R_{F_-}(s),
\]
in the extended nonnegative reals. A finite upper bound requires the pessimistic envelope actually to attain the target at finite capital. A finding that its asymptotic ceiling equals $y_0$ is not enough.

For a local error calculation suppose $F$ and $\widetilde F$ differ by at most $\epsilon$ on an interval containing both target crossings, both are increasing there, and $\partial F/\partial k\geq m>0$. If both exact crossings exist, then $|k^*_{F}-k^*_{\widetilde F}|\leq\epsilon/m$: at the second crossing the first technology is within $\epsilon$ of its target, and the mean-value bound converts output error to capital error. Near an ecological ceiling, $m$ can be arbitrarily small, so accurate output prediction need not imply accurate replacement cost. With an active baseline constraint $k\geq k_0$, apply the same argument to unconstrained crossings and the nonexpansive map $k\mapsto\max\{k,k_0\}$ when these crossings exist.

\section{Status and checked source}
The CES frontier, concave observational-equivalence example and envelope bounds are the achieved results. Identification of the actual production function, ecological dynamics, feasible investment set and service aggregation remains unresolved. None of the examples estimates a real basin's threshold.

Stefano Giglio, Theresa Kuchler, Johannes Stroebel and Olivier Wang (2024), \emph{The Economics of Biodiversity Loss}, NBER WP 32678, primary abstract inspected for the distinction between substitutable species and essential ecosystem functions:
\url{https://www.nber.org/papers/w32678}. The replacement-cost formulas and counterexample above are separate calculations, not attributed to the paper.

\end{document}
